\section{Crear y probar reglas CQL}

\subsection{Crear una regla}

Entrar a \texttt{Reglas CQL} y presionar \texttt{Nueva regla}. El sistema abre
un editor con CQL inicial. Completar la metadata:

\begin{itemize}
  \item \textbf{Titulo}: nombre visible para la clase.
  \item \textbf{Nombre CQL}: nombre de la libreria CQL.
  \item \textbf{Version}: por ejemplo \texttt{0.1.0}.
  \item \textbf{Hook}: normalmente \texttt{patient-view}.
  \item \textbf{Expresion booleana}: por ejemplo \texttt{Aplica}.
  \item \textbf{Summary}: texto corto de la card.
  \item \textbf{Detail}: explicacion educativa.
  \item \textbf{Indicator}: \texttt{info}, \texttt{warning} o \texttt{critical}.
\end{itemize}

\rceScreenshot{07-reglas-catalogo.png}{Catalogo de reglas CQL.}
\rceScreenshot{08-editor-cql.png}{Editor CQL integrado dentro del RCE.}

\subsection{Validar y ver ELM}

Presionar \texttt{Guardar} y luego \texttt{Validar}. Si el CQL es valido, el
backend envia el texto al traductor y muestra el ELM generado. Si hay errores,
se muestran diagnosticos por linea y columna.

\rceScreenshot{09-validacion-elm.png}{Vista de ELM o diagnosticos despues de validar una regla.}

\subsection{Regla simple por edad}

Esta regla se activa para pacientes adultos:

\begin{lstlisting}
library RceAdultPatient version '0.1.0'

using FHIR version '4.0.1'

context Patient

define "Aplica":
  AgeInYears() >= 18
\end{lstlisting}

Probar primero con un paciente menor de 18 anos para ver \texttt{No aplica}.
Luego cambiar la fecha de nacimiento desde la ficha del paciente para que tenga
18 anos o mas y reevaluar.

\rceScreenshot{10-prueba-no-aplica.png}{Prueba de regla con resultado No aplica antes de modificar datos.}

\subsection{Regla de diabetes educativa por HbA1c}

Esta regla muestra una card critica si el paciente adulto tiene HbA1c mayor o
igual a 6.5. Es solo una regla docente, no un diagnostico clinico real.

Metadata sugerida:

\begin{lstlisting}
Titulo: Sospecha de diabetes por HbA1c
Nombre CQL: RceDiabetesRiskHba1c
Version: 0.1.0
Hook: patient-view
Expresion booleana: Aplica
Summary: Sospecha de diabetes por HbA1c
Detail: El paciente adulto tiene HbA1c mayor o igual a 6.5% en el laboratorio sintetico del sandbox.
Indicator: critical
\end{lstlisting}

CQL:

\begin{lstlisting}
library RceDiabetesRiskHba1c version '0.1.0'

using FHIR version '4.0.1'

codesystem "LOINC": 'http://loinc.org'

code "Hemoglobin A1c": '4548-4' from "LOINC" display 'Hemoglobin A1c/Hemoglobin.total in Blood'

context Patient

define "Aplica":
  AgeInYears() >= 18
    and exists (
      [Observation: "Hemoglobin A1c"] O
        where O.status.value = 'final'
          and (O.value as FHIR.Quantity).value.value >= 6.5
    )
\end{lstlisting}

Demo recomendada:

\begin{enumerate}
  \item Elegir un paciente adulto.
  \item En \texttt{Editar dato}, dejar HbA1c en \texttt{5.6}.
  \item Guardar y reevaluar. La regla no deberia generar card.
  \item Cambiar HbA1c a \texttt{7.2}.
  \item Guardar cambios.
  \item Confirmar que aparece la card \texttt{Sospecha de diabetes por HbA1c}.
\end{enumerate}

\rceScreenshot{11-card-activa.png}{Card CDS visible despues de modificar un dato clinico.}
